\bigskip\Needspace{10\baselineskip}
\pdfbookmark[1]{Further Report. Ceres. February 1802}{ceres180202}
\begin{center}
{\small FURTHER REPORT ON THE NEW MAJOR PLANET.\par}
[CERES.]\par
\medskip\hrule\medskip
Monatliche Correspondenz. February 1802.\par
\medskip\hrule
\end{center}\medskip
—\quad —\quad —\quad “The agreement of the declinations in \textsc{Piazzi}'s observations, both with one another and with the calculation, really is admirable, and it is precisely this that strengthens my hope that my elements will be accurate enough for us to find it. For it is precisely the latitude observations, which depend most on the declinations, that have the greatest influence on the agreement of the elements. A change of $20''$ in latitude, if applied positively to the outer observations and negatively to the middle ones, or the other way around, would change the ellipse completely. Another remark goes with this one. It will gain its full value only when the recovery has succeeded and confirmed the elements derived from the observations. In that event we will owe the recovery chiefly to the large inclination of the orbit. If the orbit coincided with the ecliptic, we would probably have to give up on the recovery. At any rate, after a year the uncertainty in longitude would be extraordinarily large and might amount to a couple of signs, although this zone would then have hardly any width at all. It is easy to see why. In that case four longitudes would be needed to determine the ellipse, so that two outer observations and one middle observation could be satisfied perfectly by infinitely many ellipses, and perhaps also by parabolas and hyperbolas. Given the short series of observations, it is easy to see
\SourceLeaf{210}%
that elements representing two outer observations and one middle observation can also depart little from the whole series. Thus the conic section would be almost entirely undetermined...”

Now that \textit{Ceres} has been recovered, it is evident how sound Dr.\ \textsc{Gauss}'s judgment was in this matter, and the close agreement of his elements with the present observations confirms it. Indeed, without Dr.\ \textsc{Gauss}'s ingenious efforts and calculations we might not have found \textit{Ceres} again. The greater and finer share of the credit therefore belongs to him. The worthy Dr.\ \textsc{Olbers} is also of this opinion. He writes to the editor, “You will have noticed with pleasure how closely Dr.\ \textsc{Gauss}'s ellipse agrees with the observations of \textit{Ceres}. Do tell this worthy scholar, with an expression of my very special esteem. Without his laborious investigations into the elliptical elements of this planet, we might not have recovered it at all. I, at least, would not have looked for it so far to the east.”
\begin{flushright}[\textsc{von Zach}.]\end{flushright}

